 $\star$ EXERCISE. (a) Suppose $A$ and $B$ are objects in a category $\mathcal{C}$. Give a bijection between the natural transformations $h_A \to h_B$ of covariant functors $\mathcal{C} \to \mathrm{Sets}$ (see Example 1.1.14 for the definition) and the morphisms $B \to A$. (b) State and prove the corresponding fact for contravariant functors $h_A$ (see Example 1.1.20). Remark: A contravariant functor $F$ from $\mathcal{C}$ to Sets is said to be representable if there is a natural isomorphism $\xi : F \to h_A$. Thus the representing object $A$ is determined up to unique isomorphism by the pair $(F,\xi)$. There is a similar definition for covariant functors. (We will revisit this in §7.6, and this problem will appear again as Exercise 7.6.C. The element $\xi^{-1}(\operatorname{id}_A) \in F(A)$ is often called the “universal object”; do you see why?) (c) Yoneda’s Lemma. Suppose $F$ is a covariant functor $\mathcal{C} \to \mathrm{Sets}$, and $A \in \mathcal{C}$. Give a bijection between the natural transformations $h_A \to F$ and $F(A)$. (The corresponding fact for contravariant functors is essentially Exercise 10.1.B.)
